\subsection{Essentially Unique Solvability and Simple iSCMs}

Whereas the existence of solution functions is important, their uniqueness also plays a pivotal role in the theory of (cyclic) iSCMs.
A technical subtlety is that like existence, we will only demand uniqueness up to certain null sets. 
\begin{Def}\label{def:scm_unique_solvability}
  Let $M = \lt \Sin, \Send, \Srnd, \CX, P, f \rt$ be an iSCM.
  We say that $M$ is \emph{essentially uniquely solvable} if there exists a measurable function $g : \CX_J \times \CX_W \to \CX_V$ that satisfies
    $$x_V = f(x) \iff x_V = g(x_J,x_W)\quad M\as.$$
  For a subset $L \ins V$, we say that $M$ is \emph{essentially uniquely solvable w.r.t.\ $L$} if there exists a measurable function $g^{[L]} : \CX_J \times \CX_{(V\sm L)} \times \CX_W \to \CX_L$ that satisfies
    $$x_L = f_L(x) \iff x_L = g^{[L]}(x_J,x_{(V\sm L)},x_W)\quad M\as.$$
\end{Def}
Note that such a function $g$ is a solution function of $M$, and $g^{[L]}$ is a partial solution function of $M$ w.r.t.\ $L$.
Hence, if $M$ is essentially uniquely solvable (w.r.t.\ $L$), then $M$ is solvable (w.r.t.\ $L$).
In addition to the existence of (partial) uniqueness functions, their uniqueness up to certain null sets follows.
\begin{Def}\label{def:scm_essentially_unique_solution_function}
  Let $M = \lt \Sin, \Send, \Srnd, \CX, P, f \rt$ be an iSCM.
  We say that its solution function is \emph{essentially unique} if $g = \tilde{g}\ M\as$ for any two solution functions $g,\tilde{g}$ of $M$.
  We say that its partial solution function w.r.t.\ $L \ins V$ is \emph{essentially unique} if $g^{[L]} = \tilde{g}^{[L]}\ M\as$ for any two partial solution functions $g^{[L]}, \tilde{g}^{[L]}$ of $M$ w.r.t.\ $L$.
\end{Def}
The following observation is a simple consequence of the definitions that we will employ often.
\begin{Prp}\label{prp:scm_essentially_unique_solution_functions}
  Let $M = \lt \Sin, \Send, \Srnd, \CX, P, f \rt$ be an iSCM.
  Suppose that $M$ is essentially uniquely solvable w.r.t.\ $L \ins V$.
  Then for any (partial) solution function $\tilde{g}^{[L]}$ of $M$ w.r.t.\ $L$ we have:
    $$x_L = f_L(x) \iff x_L = \tilde{g}^{[L]}(x_J,x_{(V\sm L)},x_W)\quad M\as.$$
  Hence, the (partial) solution function of $M$ w.r.t.\ $L$ is essentially unique.
\end{Prp}
\begin{proof}
  Since $M$ is essentially uniquely solvable w.r.t.\ $L$, there exists a measurable function $g^{[L]} : \CX_J \times \CX_{(V\sm L)} \times \CX_W \to \CX_L$ that satisfies
    $$x_L = f_L(x) \iff x_L = g^{[L]}(x_J,x_{(V\sm L)},x_W)\quad M\as.$$
  Then
    $$[\forall x_V \in \CX_V : (x_L = f_L(x_J,x_{(V\sm L)},x_W,x_L) \iff x_L = g^{[L]}(x_J,x_{(V\sm L)},x_W))]\quad M\as.$$
  Let $\tilde{g}^{[L]} : \CX_J \times \CX_{(V\sm L)} \times \CX_W \to \CX_L$ be a partial solution function of $M$ w.r.t.\ $L$.
  Then:
      $$\tilde{g}^{[L]}(x_J, x_{(V\sm L)}, x_W) = f_L(x_J, x_{(V\sm L)}, x_W, \tilde{g}^{[L]}(x_J, x_{(V\sm L)}, x_W)) \quad M\as.$$
  Combining these we conclude that
    $$g^{[L]}(x_J,x_{(V\sm L)},x_W) = \tilde{g}^{[L]}(x_J,x_{(V\sm L)},x_W) \quad M\as.$$
  Together with the property of $g^{[L]}$, this also implies:
    $$x_L = f_L(x) \iff x_L = \tilde{g}^{[L]}(x_J,x_{(V\sm L)},x_W)\quad M\as.$$
\end{proof}
Essentially uniquely solvable iSCMs not only have essentially unique solution functions, but also unique distributions of potential outcomes and unique Markov kernels of solutions.
\begin{Thm}\label{thm:scm_essentially_unique_solvability_implies}
  Let $M = \lt \Sin, \Send, \Srnd, \CX, P, f \rt$ be an iSCM.
  If $M$ is essentially uniquely solvable, then:
  \begin{enumerate}
    \item Its solution function is essentially unique;
    \item The distribution of a potential outcome $X^{\doit(\xJ)}_{V\cup W}$ for input $\xJ\in\XJ$ is unique and given by 
      $$X_{V\cup W}^{\doit(\xJ)} \sim (g^{\doit(\xJ)},\id_{\CXW})_* P_M(X_W)$$
      where $g^{\doit(\xJ)} : \CXW \to \CXV : \xW \mapsto g(x_J,\xW)$ with $g : \CXJ \times \CXW \to \CXV$ a solution function of $M$.
    \item $M$ has a unique Markov kernel given by
      $$P_{M}(\XV,\XW \mid\doit(\XJ)) = (g,\id_{\CXW})_* \lp \bigg( \bigotimes_{w\in W} P_w(X_w) \bigg) \otimes \delta(X_J|X_J) \rp$$
      for $g : \CXJ \times \CXW \to \CXV$ a solution function of $M$.
  \end{enumerate}
\end{Thm}
\begin{proof}
  If $M$ is essentially uniquely solvable, it clearly has a solution function.
  By Proposition~\ref{prp:using_solution_functions}, $M$ then also has potential outcomes and solutions.
  We proceed to show their (essential) uniqueness.

%  for all x: f(x) iff g(x)
%  implies
%  for all x: f(x) iff for all x: g(x)

  \begin{enumerate}
\begin{comment}
  \item This is just the special case $L = V$ of Proposition~\ref{prp:scm_essentially_unique_solution_functions}.
  \item Let $g$ be a measurable function that satisfies
    $$x_V = f(x) \iff x_V = g(x_J,x_W)\quad M\as.$$
    Then
      $$[\forall x_V \in \CX_V : (x_V = f(x_J,x_W,x_V) \iff x_V = g(x_J,x_W))]\quad M\as.$$
    Let $\tilde{g}$ be a solution function of $M$.
    Then
      $$\tilde{g}(x_J, x_W) = f(x_J, x_W, \tilde{g}(x_J, x_W)) \quad M\as.$$
    Combining these we conclude that
      $$g(x_J,x_W) = \tilde{g}(x_J,x_W) \quad M\as.$$
    Together with the property of $g$, it implies:
      $$x_V = f(x) \iff x_V = \tilde{g}(x_J,x_W)\quad M\as.$$
      
    \item
%  Let $g, \tilde{g} : \CX_J \times \CX_W \to \CX_V$ be two solution functions of $M$.
%  Let $N = \{(x_J,x_W) \in \CX_J \times \CX_W : g(x_J,x_W) \ne f(x_J,x_W,g(x_J,x_W))\}$ be the set of exogenous states for which $g$ fails to solve the structural equations.
%  Let $\tilde N = \{(x_J,x_W) \in \CX_J \times \CX_W : \tilde g(x_J,x_W) \ne f(x_J,x_W,\tilde g(x_J,x_W))\}$ be the set of exogenous states for which $\tilde g$ fails to solve the structural equations.
%  Then $N$, $\tilde N$ and $U^c$ are $M$-null sets.
%  Hence $N \cup \tilde N \cup U^c$ is an $M$-null set.
%  On its complement, $N^c \cap \tilde N^c \cap U$, we have $g = \tilde{g}$ because of the property of $U$.
%  Hence $g = \tilde g\ M\as$.
%  \Joris{Note that for this property, it suffices if $U^c$ is an $M$-null set, not necessarily contained in a measurable $M$-null set.}
      This follows from 1.
    Let $g, \tilde{g} : \CX_J \times \CX_W \to \CX_V$ be two solution functions of $M$.
    Then
      $$x_V = f(x) \iff x_V = g(x_J,x_W)\quad M\as.$$
      $$x_V = f(x) \iff x_V = \tilde{g}(x_J,x_W)\quad M\as.$$
    Hence
      $$x_V = g(x_J,x_W) \iff x_V = \tilde{g}(x_J,x_W)\quad M\as.$$
    Hence
      $$g(x_J,x_W) = \tilde{g}(x_J,x_W) \quad M\as.$$
\end{comment}

    \item
    Let $x_J \in \CX_J$ and $X^{\doit(x_J)}_{V,W}$ a potential outcome of $M$ with input $x_J$.
    Let $g$ be a solution function of $M$.
    Then:
    $$[x_V = f(x_J,x_W,x_V) \iff x_V = g(x_J,x_W)] \quad M\as.$$
    With Lemma~\ref{lem:M-as_to_rv-as} we conclude:
    $$[X^{\doit(x_J)}_V = f(x_J,X^{\doit(x_J)}_W,X^{\doit(x_J)}_V) \iff X^{\doit(x_J)}_V = g(x_J,X^{\doit(x_J)}_W)] \quad \text{a.s.}$$
    This implies:
    $$[X^{\doit(x_J)}_V = f(x_J,X^{\doit(x_J)}_W,X^{\doit(x_J)}) \text{ a.s.}] \iff [X^{\doit(x_J)}_V = g(x_J,X^{\doit(x_J)}_W) \text{ a.s.}].$$
    Since the l.h.s.\ holds by assumption, we conclude that
    $$X_{V\cup W}^{\doit(x_J)} = (g(x_J,X^{\doit(x_J)}_W), X^{\doit(x_J)}_W)\ \text{a.s.},$$
    which in turn implies 
    $$X_{V\cup W}^{\doit(\xJ)} \sim (g^{\doit(\xJ)},\id_{\CXW})_* P(X^{\doit(x_J)}_W)$$
    and therefore
    $$X_{V\cup W}^{\doit(\xJ)} \sim (g^{\doit(\xJ)},\id_{\CXW})_* P_M(X_W).$$

%  Let $\tilde g$ be a solution function of $M$ such that $\tilde g(x_J,x_W) = f(x_J,x_W,\tilde g(x_J,x_W))$ on $U$
%  (its existence is shown in the proof of Proposition~\ref{thm:scm_solvability_sufficient_conditions}).
%
%  By assumption, $U_{x_J}^c = \{x_W \in \CX_W : (x_J,x_W) \not\in U\}$ is a $P(X_W)$-null set.
%  That is, there exists a set $N \supseteq U_{x_J}^c$ with $P(X_W \in N) = 0$.
%  Hence, $P(X_W^{\doit(x_J)} \in N) = 0$, that is, $X_W^{\doit(x_J)} \in U_{x_J}$ a.s..
%  Furthermore, by assumption $(X_W^{\doit(x_J)},X_V^{\doit(x_J)}) \in S_{x_J}$ a.s.,
%  where $S_{x_J} := \{x_{W\cup V} \in \CX_{W\cup V} : x_V = f(x_J,x_W,x_V)\}$.
%  Combining these, we conclude $X_V^{\doit(x_J)} = \tilde g(x_J,X^{\doit(x_J)}_W)$ a.s..
%  This obviously implies that $X_{V\cup W}^{\doit(x_J)} = (\tilde g(x_J,X^{\doit(x_J)}_W), X^{\doit(x_J)}_W)$ a.s.,
%  which then implies $X_{V\cup W}^{\doit(\xJ)} \sim (\tilde g^{\doit(\xJ)},\id_{\CXW})_* P(X_W)$.
%  Finally, note that since $g = \tilde g\ M\as$, for each $x_J \in \CX_J$: $\tilde g^{\doit(\xJ)} = g^{\doit(x_J)}\ P(X_W)\as$.
%  Hence
%    $$X_{V\cup W}^{\doit(\xJ)} \sim (g^{\doit(\xJ)},\id_{\CXW})_* P(X_W)$$

  \item
  Let $K(\XV,\XW \mid \XJ)$ denote the Markov kernel for $M$ corresponding to a solution $X : \C{U} \times \CXJ \to \CX$.
  For any $x_J \in \CX_J$, the random variable $X^{\doit(x_J)} : \C{U} \to \CX : u \mapsto X(u,x_J)$ is a potential outcome of $M$.
  Its distribution coincides with that of $K(\XV,\XW \mid \XJ=\xJ)$, by definition.
  The claim now follows from 3.
  \end{enumerate}
\end{proof}
For the special case of no exogenous input variables ($\Sin = \emptyset$), this means that essentially uniquely solvable SCMs induce a unique distribution.

If an iSCM is essentially uniquely solvable, this does not necessarily mean that it is still essentially uniquely solvable after performing some intervention. 
To avoid the complications introduced in case solutions are absent, or present but not unique, we will henceforth make strong assumptions regarding the existence and uniqueness of solutions. 
We will (mostly) restrict our attention to a subclass of iSCMs that we refer to as \emph{simple iSCMs}, which are iSCMs that are essentially uniquely solvable and remain so after any hard intervention:
\begin{Def}\label{def:simple_scm}
  %An iSCM $M = \lt \Sin, \Send, \Srnd, \CX, P, f \rt$ is called \emph{simple} if the intervened iSCM $M_{\doit(T)}$ is essentially uniquely solvable for all $T \subseteq \Send$. %\cup \Sin$. % \cup \Srnd$.
  %\Joris{It seems that this leads to problems with marginalization!}
  An iSCM $M = \lt \Sin, \Send, \Srnd, \CX, P, f \rt$ is called \emph{simple} if for all $T \subseteq \Send$, $M$ is essentially uniquely solvable w.r.t.\ $T$.
\end{Def}
Note that this includes essentially unique solvability of $M$ itself for $T = \emptyset$.

\begin{Rem}\label{rem:scm_simple_sufficient_condition}
An equivalent condition for $M$ to be essentially uniquely solvable w.r.t.\ $L$ is provided by Proposition~\ref{prp:scm_partial_solvability}.
Thereby, $M$ is a simple iSCM if and only if for each $L \ins V$ the following holds:
$$[\exists! x_L \in \CX_L : x_L = f_L(x_J,x_{V\sm L},x_W,x_L)]$$
up to a measurable $M$-null set \TODOmas.

For $M$ an SCM without input nodes (that is, $J = \emptyset$), every $M$-null set is contained in a measurable $M$-null set.
Hence, in that case Proposition~\ref{prp:scm_partial_solvability} implies that $M$ is essentially uniquely solvable if and only if
$$\exists! x_V \in \CX_V : x_V = f(x_V,x_W)\quad P\as.$$
\end{Rem}
%\begin{proof}
%  $\impliedby$: According to Proposition~\ref{prp:scm_partial_solvability}:
%  If 
%  $$\exists! x_V \in \CX_V : x_V = f(x_V,x_W)\quad P\as$$
%  then there exists a measurable function $g: \CXW \to \CX_V$ such that
%    $$x_V = f(x_V,x_W) \iff x_V = g(x_W)\quad P\as.$$
%
%  $\implies$: if there exists a measurable function $g : \CX_W \to \CX_V$ that satisfies
%    $$x_V = f(x_V,x_W) \iff x_V = g(x_W)\quad P\as,$$
%  then clearly
%  $$\exists! x_V \in \CX_V : x_V = f(x_V,x_W)\quad P\as.$$
%\end{proof}

\begin{Eg}\label{eg:simple_scm}
Consider an iSCM with structural equations
  \begin{align*}
    X_1 &= W_1 \\
    X_2 &= W_2 \\
    X_3 &= X_1 X_4 + W_3 \\
    X_4 &= X_2 X_3 + W_4 
  \end{align*}
where the $X$'s are considered real-valued endogenous variables and the $W$'s exogenous variables with domains $(-1,1) \subset \RN$.
We can solve the system of structural equations for $X$ in terms of $W$:
  \begin{align*}
    X_1 &= W_1 \\
    X_2 &= W_2 \\
    X_3 &= \frac{W_3 + W_1 W_4}{1 - W_1 W_2} \\
    X_4 &= \frac{W_2 W_3 + W_4}{1 - W_2 W_1}
  \end{align*}
which gives a unique solution function.
Similarly, we can take any subset of the structural equations and solve it for the variables appearing on the l.h.s.\ of the equations in the subset, and obtain a unique solution function. 
For example, only solving the structural equations for $X_3$ and $X_4$, we obtain:
  \begin{align*}
    X_3 &= \frac{W_3 + X_1 W_4}{1 - X_1 X_2} \\
    X_4 &= \frac{X_2 W_3 + W_4}{1 - X_2 X_1}
  \end{align*}
where the variables $X_1$ and $X_2$ are now considered as exogenous input variables (instead of endogenous variables).
Hence, any subset of the structural equations has a unique solution for the variables appearing on the l.h.s.\ in terms of the remaining ones on the r.h.s., which means that this iSCM is simple.

If we extend the range of possible values of the exogenous random variables to $\CX_W = \RN^2$, then we may still end up with essentially unique solution functions as long as $\{(w_1,w_2) \in \RN^2 : 1 = w_1 w_2\}$ is an $M$-null set (which depends on the distribution $P(X_W)$ we would pick).
\end{Eg}

One of the very convenient properties of simple iSCMs is that after any hard intervention, they still have a uniquely defined Markov kernel.
Simplicity is preserved by hard interventions and by adding intervention variables.
\begin{Prp}\label{prp:scm_simplicity_preserved_by_interventions}
  Let $M = \lt \Sin, \Send, \Srnd, \CX, P, f \rt$ be a simple iSCM. Then:
    \begin{enumerate}
      \item $M_{\doit(T\dots)}$ is simple for $T \subseteq V \cup J$ (hard interventions of all three types);
      \item $M_{\doit(I_B)}$ is simple for $B \subseteq V \cup J$ (adding intervention variables).
    \end{enumerate}
\end{Prp}
\begin{proof}
  Because $M$ is simple, for all $T \ins V$ there exists a measurable function $g_{\doit(T)} : \CX_{J \cup T} \times \CX_W \to \CX_{V\sm T}$ such that:
    $$x_{V\sm T} = f_{V\sm T}(x) \iff x_{V\sm T} = g_{\doit(T)}(x_J, x_T, x_W)\quad M\as.$$
  Let $\tilde{M} = \langle \tilde{J}, \tilde{V}, \tilde{W}, \CX, \tilde{P}, \tilde{f}\rangle$ denote an intervened iSCM (we will consider special cases below).
  To prove that $\tilde{M}$ is simple, it suffices to show that for all $\tilde{T} \ins \tilde{V}$ there exists a measurable function
    $\tilde{g}_{\doit(\tilde{T})} : \CX_{\tilde{J} \cup \tilde{T}} \times \CX_{\tilde{W}} \to \CX_{\tilde{V}\sm \tilde{T}}$ such that:
    $$x_{\tilde{V}\sm \tilde{T}} = \tilde{f}_{\tilde{V}\sm \tilde{T}}(x) \iff x_{\tilde{V}\sm \tilde{T}} = \tilde{g}_{\doit(\tilde{T})}(x_{\tilde{J}}, x_{\tilde{T}}, x_{\tilde{W}})\quad \tilde{M}\as,$$
  as long as $M\as$ implies $\tilde{M}\as$.

  \begin{enumerate}
    \item 
      \begin{enumerate}[label=(\roman*)]
        \item
      Let $T \ins J \cup V$. 
      Consider $\tilde{M} = M_{\doit(T)}$, that is, $\tilde{J} = J \cup T$, $\tilde{V} = V \sm T$, $\tilde{W} = W$, $\tilde{P} = P$, $\tilde{f} = f_{\tilde{V}}$.
      Since $\tilde{M}_{\doit(\tilde{T})} = M_{\doit(T \cup \tilde{T})}$ for $\tilde{T} \ins \tilde{V}$, we can just take $\tilde{g}_{\doit(\tilde{T})} = g_{\doit((T \sm J) \cup \tilde{T})}$.
      Note that indeed $M\as$ implies $M_{\doit(T)}\as$.
      Hence $M_{\doit(T)}$ is simple.
    \item
      Let $T \ins V \cup J$, and $\xi_T \in \CX_T$. 
      Consider $\tilde{M} = M_{\doit(x_T=\xi_T)}$, that is, $\tilde{J} = J \sm T$, $\tilde{V} = V \cup T$, $\tilde{W} = W$, $\tilde{P} = P$, $\tilde{f} = (f_{V\sm T},\xi_T)$. 
      Since for $\tilde{T} \ins \tilde{V}$:
      $$\tilde{M}_{\doit(\tilde{T})} = (M_{\doit(x_{T\sm \tilde{T}}=\xi_{T\sm \tilde{T}})})_{\doit(\tilde{T})},$$
      as solution function $\tilde{g}_{\doit(\tilde{T})}$ we can take the one for $M_{\doit(T \cup \tilde{T})}$, evaluated in $x_{T\sm \tilde{T}}=\xi_{T\sm \tilde{T}}$, and with additional output $T \sm \tilde{T}$ with value $\xi_{T\sm \tilde{T}}$:
      $$\tilde{g}_{\doit(\tilde{T})}(x_{\tilde{J}}, x_{\tilde{T}}, x_W) = \big(g_{\doit((T\sm J) \cup \tilde{T})}(\xi_{T \sm \tilde{T}},x_{\tilde{T} \cup J},x_W),\xi_{T \sm \tilde{T}}\big).$$
      Note that indeed $M\as$ implies $M_{\doit(x_T=\xi_T)}\as$.
      Hence $M_{\doit(x_T=\xi_T)}$ is simple.
    \item
      Let $T \ins V \cup J$ and $Q_T \in \bigotimes_{t\in T} \C{P}(\C{X}_t)$. 
      Consider $\tilde{M} = M_{\doit(\XT\sim Q_T)}$, that is, $\tilde{J} = J\sm T$, $\tilde{V} = V\sm T$, $\tilde{W} = W\cup T$, $\tilde{P}=P \otimes Q_T$, $\tilde{f}=f_{V\sm T}$.
      Since for $\tilde{T} \ins \tilde{V}$:
      $$\tilde{M}_{\doit(\tilde{T})} = (M_{\doit(T \cup \tilde{T})})_{\doit(\XT\sim Q_T)}$$
      we can take
      $\tilde{g}_{\doit(\tilde{T})} = g_{\doit((T \sm J) \cup \tilde{T})}$.
      Note that indeed $M\as$ implies $M_{\doit(\XT\sim Q_T)}\as$.
      Hence $M_{\doit(\XT\sim Q_T)}$ is simple.
      \end{enumerate}
    \item
      For $B \ins J$, the operation of adding intervention variables $\doit(I_{B})$ only adds indices $I_j := j$ for $j \in B \cap J$.
      Therefore, we can consider $B \subseteq V$ without loss of generality.
      Furthermore, by using induction we only need to consider the case $B = \{b\} \ins V$ by means of Proposition~\ref{prp:adding_intervention_variables_one_by_one}.
      Hence consider $\tilde{M} = M_{\doit(I_b)}$, that is, $\tilde J = J \cup \{I_b\}$, $\tilde V = V$, $\tilde W = W$, $\tilde \CX = \CX \times \CX_{I_b}$ with $\CX_{I_b} := \CX_b \,\dot\cup\, \{\star\}$, $\tilde P = P$, and $\tilde f$ given by
  \begin{align*}
    & \tilde{f}_b(x_{I_b},x_{V \cup J \cup W}) := \begin{cases}
      f_b(x_{V \cup J \cup W}) & x_{I_b} = \star \\
      x_{I_b} & x_{I_b} \in \CX_b
    \end{cases}\\
    v \in V \sm \{b\}: \qquad & \tilde{f}_v(x_{I_b},x_{V \cup J \cup W}) := f_v(x_{V \cup J \cup W})
  \end{align*}
     % For $b \in B$, $x_{I_b} = \star$ encodes that there is no intervention on $b$, while $x_{I_b} \ne \star$ encodes that the hard intervention $\doit(X_b = x_{I_b})$ is performed.
    Let $\tilde{T} \ins \tilde{V}$.
    Consider first the case $b \notin \tilde{T}$. Then:
      $$\tilde{M}_{\doit(\tilde{T})} = (M_{\doit(\tilde{T})})_{\doit(I_b)}.$$
    We can take:
      $$\tilde{g}_{\doit(\tilde{T})}(x_{I_b},x_J,x_{\tilde{T}},x_W) = \begin{cases}
        g_{\doit(\tilde{T})}(x_J,x_{\tilde{T}},x_W) & x_{I_b} = \star \\
        \big(g_{\doit(\tilde{T} \cup \{b\})}(x_J,x_{\tilde{T}},x_b=x_{I_b},x_W),x_{I_b}\big) & x_{I_b} \in \CX_b 
      \end{cases}$$
    as one can check that 
      $$
      \begin{cases}
      x_{V\sm \tilde{T}} = f_{V\sm \tilde{T}}(x_{J \cup W \cup V}) \iff x_{V\sm \tilde{T}} = g_{\doit(\tilde{T})}(x_J, x_{\tilde{T}}, x_W)\quad M\as. \\
      x_{V\sm (\tilde{T} \cup \{b\})} = f_{V\sm (\tilde{T}\cup\{b\})}(x_{J \cup W \cup V}) \iff x_{V\sm (\tilde{T}\cup\{b\})} = g_{\doit(\tilde{T}\cup\{b\})}(x_J, x_{\tilde{T}\cup\{b\}}, x_W)\quad M\as.
      \end{cases}$$
    implies that for each $x_{I_b} \in \CX_{I_b}$:
      $$x_{V\sm \tilde{T}} = \tilde{f}_{V\sm \tilde{T}}(x_{I_b},x_{J\cup W\cup V}) \iff x_{V\sm \tilde{T}} = \tilde{g}_{\doit(\tilde{T})}(x_{I_b}, x_J, x_{\tilde{T}}, x_W)\quad M\as.$$

    Now consider the case $b \in \tilde{T}$. Then:
      $$\tilde{M}_{\doit(\tilde{T})} = ((M_{\doit(\tilde{T}\sm \{b\})})_{\doit(I_b)})_{\doit(b)}.$$
    This is just $M_{\doit(\tilde{T})}$ with an additional `unused' exogenous input variable $I_b$.
    We can therefore take:
      $$\tilde{g}_{\doit(\tilde{T})}(x_{I_b},x_J,x_{\tilde{T}},x_W) = g_{\doit(\tilde{T})}(x_J,x_{\tilde{T}},x_W)$$
    where one should notice that in this case the desired property holds $M_{\doit(I_b)}\as$, as it does not depend on $x_{I_b}$ and holds $M\as$.
  \end{enumerate}
  This completes the proof.
\end{proof}
Simplicity is in general not preserved by soft interventions.

\begin{Cor}\label{cor:simple_scm}
  If an iSCM $M = \lt \Sin, \Send, \Srnd, \CX, P, f \rt$ is simple, then its Markov kernel $P_{M}(X_V, X_W \mid \doit(X_J))$ and all intervened Markov kernels
%  $$P_{M}(X_{(V \cup W) \setminus T} \mid \doit(X_J),\doit(X_T)) := P_{M_{\doit(T)}}(X_{(V \cup W)\setminus T} \mid \doit(X_{J \cup T}))$$
  $$P_{M_{\doit(T)}}(X_{V \sm T},X_W \mid \doit(X_{J \cup T}))$$
  for $T \subseteq V$ % \cup W$ 
  exist and are unique.
\end{Cor}
\begin{proof}
  If $M$ is simple, then for all $T \ins V$ also $M_{\doit(T)}$ is simple (by Proposition~\ref{prp:scm_simplicity_preserved_by_interventions}) and hence essentially uniquely solvable.
  The claim now follows from Theorem~\ref{thm:scm_essentially_unique_solvability_implies}.
\end{proof}
We will make use of the following notation for intervened Markov kernels of simple iSCMs.
\begin{Not}\label{not:scm_do_operator}
  For a simple iSCM $M$, we write for $T \subseteq V$:
  $$P_{M}(X_{V \sm T}, X_W \mid \doit(X_J),\doit(X_T)) := P_{M_{\doit(T)}}(X_{V \sm T},X_W \mid \doit(X_{J \cup T}))$$
%  We will refer to $P_M(X_{V\sm T} \mid \doit(X_J),\doit(X_T))$ as the \emph{interventional Markov kernels} induced by the iSCM,
%  and (for $T = \emptyset$) to $P_M(X_V \mid \doit(X_J))$ as the \emph{observable Markov kernels} induced by the iSCM. 
%  For the special case $J = T = \emptyset$, we will refer to $P_M(X_V)$ as the \emph{observational Markov kernel} of the iSCM.
  We establish the following terminology for these Markov kernels:

  \medskip
  \centerline{\begin{tabular}{lll}
    special case        & Markov kernel & name \\
    \hline
                        & $P_M(X_{V\sm T} \mid \doit(X_J),\doit(X_T))$ & \emph{interventional Markov kernels of $M$} \\
    $T = \emptyset$     & $P_M(X_V \mid \doit(X_J))$ & \emph{observable Markov kernel of $M$} \\
    $J = T = \emptyset$ & $P_M(X_V)$ & \emph{observational Markov kernel of $M$}
  \end{tabular}}
\end{Not}

The fact that these iSCMs are relatively simple to deal with (because we do not have to worry about non-existence or non-uniqueness of solutions) motivated their name.
Even though simplicity is a strong assumption, the class of simple iSCMs is more expressive than the class of L-CBNs, since it allows to model causal cycles to some extent.
In paricular, some of the examples at the beginning of this Chapter (in Section~\ref{sec:scms_motivation}) can be dealt with using the class of simple iSCMs.

\begin{Exc}
  Show that Example~\ref{eg:coarse_grained_cbn} provides no difficulties when modeling using simple iSCMs.
  Is the class of iSCMs generally closed under merging variables?
\end{Exc}

\begin{Exc}
  Show that the mixture model in Example~\ref{eg:ohms_law} can be modeled as a simple iSCM.
  Is the class of iSCMs generally closed under taking mixtures of iSCMs?
\end{Exc}

\subsubsection{Linear iSCMs}

For iSCMs whose causal mechanism is linear in terms of the endogenous variables (up to a certain null set), we can give a sufficient condition for essentially unique solvability, and explicitly write down the form of their solution function.
\begin{Def}\label{def:linear_scms}
Let $\mathbb{K}$ be a field (e.g., $\mathbb{K} = \RN$ or $\mathbb{K} = \mathbb{C}$).
Let $M=\lt \Sin, \Send, \Srnd, \CX, P, f \rt$ be an iSCM, where each $\CX_v$ is a finite-dimensional Polish vector space over $\mathbb{K}$, and let $\CX_A = \bigoplus_{v \in A} \CX_a$ for $A \ins V$.
Then $M$ is called \emph{essentially linear} if each component of the causal mechanism is (up to an $M$-null set) an affine combination of endogenous variables with coefficients that may be functions of the exogenous variables, i.e., of $f$ is of the form
$$
  f(x_J,x_W,x_V) = B(x_J,x_W) x_V + c(x_J,x_W) \qquad M\as,
$$
where for all $x_J \in \CX_J$ and all $x_W \in \CX_W$, $B(x_J,x_W) : \CX_V \to \CX_V$ is linear and $c(x_J,x_W) \in \CX_V$.
Since $\CX_V = \bigoplus_{v \in V} \CX_v$, we can also write this as
$$
  f_v(x_J,x_W,x_V) = \sum_{u\in V} B_{vu}(x_J,x_W) x_u + c_v(x_J,x_W) \qquad M\as,
$$
where for all $x_J \in \CX_J$ and all $x_W \in \CX_W$, $B_{vu}(x_J,x_W) : \CX_u \to \CX_v$ is linear and $c_v(x_J,x_W) \in \CX_v$ for all $u,v \in V$.
\end{Def}
For $\mathbb{K} = \RN$ and $\CX_v = \RN$ for all $v\in V$, the mapping $B$ can be thought of as a matrix-valued measurable function $\CXJ\times\CXW \to \RN^{V\times V}$ and $c$ as a vector-valued function $\CXJ\times\CXW \to \RN^V$.

For essentially linear iSCMs, essentially unique solvability is implied by the invertability of a certain matrix.
\begin{Prp}\label{prp:linear_scms}
Let $M$ be a essentially linear iSCM as in Definition~\ref{def:linear_scms}.
%For any set $L \subseteq V$, denote by $I_L \in \RN^{L \times L}$ the identity matrix and $B_{LL}(x_J,x_W) \in \RN^{L \times L}$ the submatrix of $B(x_J,x_W)$.
%Then, the submodel $M^{[L]}$ is essentially uniquely solvable if the matrices $I_{L}-B_{LL}(x_J,x_W)$ are invertible for all $x_J \in \CX_J, x_W \in \CX_W$ up to a measurable $M^{[L]}$-null set.
%The corresponding essentially unique solution function is then:
%  \begin{equation*}\begin{split}
%    g^{[L]} &: \RN^{V\setminus L} \times \RN^{J\cup W} \to \RN^L \\
%      &: (x_{V\setminus L},x_J,x_W) \mapsto (I_L - B_{LL}(x_J,x_W))^{-1} \big(B_{L,{V\setminus L}}(x_J,x_W) x_{V\setminus L} + c_{L}(x_J,x_W)\big).
%  \end{split}\end{equation*}
Denote by $I \in \mathbb{K}^{V \times V}$ the identity matrix.
Then, $M$ is essentially uniquely solvable if the matrices $I-B(x_J,x_W)$ are invertible for all $x_J \in \CX_J, x_W \in \CX_W$ up to a measurable $M$-null set \TODOmas.
The corresponding essentially unique solution function is:
  $$g : \CX_J \times \CX_W \to \CX_V : (x_J,x_W) \mapsto (I - B(x_J,x_W))^{-1} c(x_J,x_W).$$

More generally, let $L \ins V$.
  Denote by $I_L \in \RN^{L \times L}$ the identity matrix and $B_{LL}(x_J,x_W) \in \mathbb{K}^{L \times L}$ the submatrix of $B(x_J,x_W)$.
$M$ is essentially uniquely solvable w.r.t.\ $L$ if the matrices $I_{L}-B_{LL}(x_J,x_W)$ are invertible for all $x_J \in \CX_J, x_W \in \CX_W$ up to a measurable $M$-null set \TODOmas.
The corresponding essentially unique partial solution function w.r.t.\ $L$ is:
  \begin{equation*}\begin{split}
    g^{[L]} &: \CX_J \times \CX_{V\setminus L} \times \CX_W \to \CX_L \\
      &: (x_J,x_{V\setminus L},x_W) \mapsto (I_L - B_{LL}(x_J,x_W))^{-1} \big(B_{L,{V\setminus L}}(x_J,x_W) x_{V\setminus L} + c_{L}(x_J,x_W)\big).
  \end{split}\end{equation*}
\end{Prp}
%A non-zero coefficient $B_{ij}$ for $i,j\in\C{I}$ such that $i\neq j$ corresponds with a directed edge $j\to i$ in the (augmented) graph, and a coefficient $B_{ii}=1$ for $i\in\C{I}$ corresponds with a self-cycle $i\to i$ in the (augmented) graph of the SCM. A non-zero coefficient $\Gamma_{ij}$ for $i\in\C{I}$, $j\in\C{J}$ with $\Prb_{\C{E}_j}$ a non-degenerate probability distribution over $\RN$ corresponds with a directed edge $j \to i$ in the augmented graph. A non-zero entry $(\Gamma \Gamma^T)_{ij}$ for $i,j\in\C{I}$ with $i\neq j$ such that there exists a $k\in\C{J}$ for which $\Gamma_{ik},\Gamma_{jk}\neq 0$ and $\Prb_{\C{E}_k}$ a non-degenerate probability distribution over $\RN$ corresponds with a bidirected edge $i \oto j$ in the graph of the SCM.

\subsection{Equivalence Notions}

\begin{comment}
Equivalence relations are ubiquitous in mathematics. They capture the notion that mathematical objects can be ``equivalent'' from some point of view.
\begin{Def}
  Let $Z$ be a set and $R \subseteq Z^2$ be a relation on $Z$ (i.e., a subset of ordered pairs of $Z$). $R$ is called an \emph{equivalence relation} if 
  \begin{enumerate}
    \item $R$ is reflexive: $(a,a) \in R$ for all $a \in Z$;
    \item $R$ is symmetric: $(a,b) \in R \iff (b,a) \in R$ for all $a,b \in Z$;
    \item $R$ is transitive: if $(a,b) \in R$ and $(b,c) \in R$ then $(a,c) \in R$ for all $a,b,c\in Z$.
  \end{enumerate}
\end{Def}
\end{comment}
In this section, we will discuss several important equivalence relations between iSCMs.

$M$ and $\tilde{M}$ are considered equivalent if they differ only in terms of their causal mechanisms, yet each of their structural equations has almost surely the same solutions.
\begin{Def}\label{def:equivalent_scms}
  Let $M = \lt \Sin, \Send, \Srnd, \CX, P, f \rt$ and
  $\tilde{M} = \lt \Sin, \Send, \Srnd, \CX, P, \tilde{f} \rt$ be two iSCMs that may differ only in terms of their causal mechanism. 
  We say that $M$ is \emph{equivalent} to $\tilde{M}$ and write $M \equiv \tilde{M}$ if for each $v \in \Send$,
  $$x_v = f_v(\xJ,\xV,\xW) \iff x_v = \tilde{f}_v(\xJ,\xV,\xW) \qquad M\as$$
  (note that $M\as$ means the same as $\tilde{M}\as$).
\end{Def}
Equivalent iSCMs are indeed equivalent for many purposes, and it is compatible with all properties of and operations on iSCMS that we have defined thus far.
%\begin{Lem}
%  Let $M \equiv \tilde{M}$ be equivalent iSCMs.
%  Define
%   $$S : \{(x_J,x_W,x_V) \in \CX_{J\cup W\cup V} \mid x_V = f(x_J,x_W,x_V) \},$$
%  and similarly 
%   $$\tilde{S} : \{(x_J,x_W,x_V) \in \CX_{J\cup W\cup V} \mid x_V = \tilde{f}(x_J,x_W,x_V) \}.$$
%  Then $S = \tilde{S}\ M\as$.
%\end{Lem}
%\begin{proof}
%  We want to show that this is contained in a $P(X_W)$-null set in $\CX_{J \cup W}$.
%  Since $M \equiv \tilde{M}$, $\forall x_V \in \CX_V: x_V \in S_{(x_J,x_W)} \iff x_V \in \tilde{S}_{(x_J,x_W)}$ holds $M\as$.
%  Hence $S_{(x_J,x_W)} = \tilde{S}_{(x_J,x_W)}$ holds $M\as$.
%  That means that 
%  $$\{(x_J,x_W) \in \CX_J \times \CX_W \mid S_{(x_J,x_W)} \ne \tilde{S}_{(x_J,x_W)} \}$$
%  is contained in a $M$-nullset, and hence in a $P(X_W)$-null set in $\CX_{J \cup W}$.
%\end{proof}
\begin{Prp}\label{prp:equivalent_scms}
\begin{enumerate}
  \item Equivalence is preserved by interventions:
    if $M \equiv \tilde{M}$ then:
    \begin{enumerate}
      \item $M_{\doit(T\dots)} \equiv \tilde{M}_{\doit(T\dots)}$ for $T \subseteq V \cup J$ (hard interventions of all three types);
      \item $M_{\doit(T\ot\hat{f}_T)} \equiv \tilde{M}_{\doit(T\ot\hat{f}_T)}$ for $T \subseteq V$ and $\hat f_T : \CX_J \times \CX_W \times \CX_V \to \CX_T$ a measurable function
        (soft interventions);
      \item $M_{\doit(I_B)} \equiv \tilde{M}_{\doit(I_B)}$ for $B \subseteq V \cup J$ (adding intervention variables);
    \end{enumerate}
  % \item Equivalence is preserved by twinning.
  \item Equivalent iSCMs have equivalent submodels and composition preserves equivalence;
  \item Equivalent iSCMs have the same null sets:
    if $M \equiv \tilde{M}$ then $N$ is an $M$-null set if and only if it is an $\tilde{M}$-null set, and `$M\as$' means the same as `$\tilde{M}\as$';
  \item Equivalent iSCMs have the same potential outcomes, solutions, Markov kernels, and (partial) solution functions;
  \item (Partial) solvability, essentially unique (partial) solvability, essential uniqueness of (partial) solution functions and simplicity are invariants under equivalence:
    if $M \equiv \tilde{M}$ then 
    \begin{enumerate}
      \item $M$ is solvable w.r.t.\ $L \ins V$ if and only if $\tilde{M}$ is solvable w.r.t.\ $L$,
      \item $M$ is essentially uniquely solvable w.r.t.\ $L \ins V$ if and only if it $\tilde{M}$ is essentially uniquely solvable w.r.t.\ $L$,
      \item $g^{[L]}$ is an essentially unique solution function of $M$ w.r.t.\ $L \ins V$ if and only if it is an essentially unique solution function of $\tilde{M}$ w.r.t.\ $L$,
      \item $M$ is simple if and only if $\tilde{M}$ is simple;
    \end{enumerate}
  \item Essential linearity is preserved.    
\end{enumerate}
\end{Prp}
\begin{proof}
  The first three claims follow directly from the definitions.
  
%  Let $M \equiv \tilde{M}$.
%  The condition 
%  $$\forall v \in V: x_v = f_v(\xJ,\xV,\xW) \iff x_v = \tilde{f}_v(\xJ,\xV,\xW) \qquad M\as$$
%  implies $S = \tilde{S}\ M\as$, where
%   $$S = \{(x_J,x_W,x_V) \in \CX_{J\cup W\cup V} \mid x_V = f(x_J,x_W,x_V) \},$$
%   $$\tilde{S} = \{(x_J,x_W,x_V) \in \CX_{J\cup W\cup V} \mid x_V = \tilde{f}(x_J,x_W,x_V) \}.$$
%  It even implies that $S \,\Delta\, \tilde{S}$ is measurable.
%  Indeed, for all $v \in V$, define the sets $S_v := \{x \in \CX : x_v = f_v(x)\}$ and
%  $\tilde{S}_v := \{x \in \CX : x_v = \tilde{f}_v(x)\}$, the functions
%  $$h_v : \CXJ \times \CXW \times \CXV \to \CX_v \times \CX_v : (\xJ,\xW,\xV) \mapsto (x_v, f_v(\xV,\xJ,\xW))$$
%  $$\tilde h_v : \CXJ \times \CXW \times \CXV \to \CX_v \times \CX_v : (\xJ,\xW,\xV) \mapsto (x_v, \tilde f_v(\xV,\xJ,\xW))$$
%  and the diagonals $\Delta_v = \{(x_v,x_v) : x_v \in \CX_v\} \subseteq \CX_v^2$. 
%  Since $h_v$ is measurable and $\Delta_v$ is a measurable set (because $\CX_v$ is Hausdorff), $S_v = h_v^{-1}(\Delta_v)$ is a measurable set.
%  Similarly, $\tilde S_v = \tilde h_v^{-1}(\Delta_v)$ is a measurable set.
%  Hence $S_v \,\Delta\, \tilde S_v$ is measurable.
%  The assumption $M \equiv \tilde M$ then implies that $S_v \,\Delta\, \tilde S_v$ is contained in an $M$-null set of the form $N_v = \tilde{N}_v \times \CX_V$ with $\tilde{N}_v \ins \CX_J \times \CX_W$ a measurable $M$-null set.

  We prove the fourth point.
  $M \equiv \tilde{M}$ implies that 
    \begin{equation}\label{eq:prp:equivalent_scms}
      x_V = f(x) \iff x_V = \tilde{f}(x)\ M\as.
    \end{equation}

    \begin{itemize}
      \item
  Potential outcomes are invariant.
          Let $x_J \in \CX_J$. 
          We show that if a random variable $x^{\doit(x_J)}_{V,W}$ is a potential outcome of $M$, then it is a potential outcome of $\tilde{M}$.
          Using Lemma~\ref{lem:M-as_to_rv-as}, \eref{eq:prp:equivalent_scms} implies:
                  $$X_V^{\doit(\xJ)} = f(x_J,X_W^{\doit(\xJ)},X_V^{\doit(\xJ)}) \iff X_V^{\doit(\xJ)} = \tilde{f}(x_J,X_W^{\doit(\xJ)},X_V^{\doit(\xJ)}) \text{ a.s.}.$$
          Since $X_W^{\doit(\xJ)} \sim P_M(X_W$), this proves the claim.

        \item
      Solutions are invariant.
          Let $(\C{U} \times \CXJ, K(U|X_J))$ be a transition probability space and $X : \C{U} \times \CXJ \to \CX$ be a conditional random variable.
          We show that if $X$ is a solution of $M$, then it is a solution of $\tilde{M}$.
          Using Lemma~\ref{lem:M-as_to_crv-as}, \eref{eq:prp:equivalent_scms} implies:
          $$X_V = f(X_J,X_W,X_V) \iff X_V = \tilde{f}(X_J,X_W,X_V) \text{ a.s.}.$$
          Since $K(X_W|X_J) = P_M(X_W$), this proves the claim.

        \item
    Markov kernels are invariant: this follows from the invariance of solutions.

  \item
    Partial solution functions are invariant.
      Let $L \ins V$.
      We show that if $g^{[L]}$ is a partial solution function of $M$ w.r.t.\ $L$, then it is a partial solution function of $\tilde{M}$ w.r.t\ $L$.
      By assumption, $g^{[L]} : \CXJ \times \CX_{V\sm L} \times \CXW \to \CX_L$ is a measurable function.
      Because $M\as$ and $\tilde{M}\as$ mean the same, 
        $$g^{[L]}(\xJ,x_{V\sm L},\xW) = f_L\big(\xJ, x_{V\sm L}, g^{[L]}(\xJ,x_{V\sm L},\xW), \xW\big) \quad M\as$$
      together with \eref{eq:prp:equivalent_scms} implies:
        $$g^{[L]}(\xJ,x_{V\sm L},\xW) = \tilde{f}_L\big(\xJ, x_{V\sm L}, g^{[L]}(\xJ,x_{V\sm L},\xW), \xW\big) \quad \tilde{M}\as,$$
        which proves the claim.
  \end{itemize}

  The fifth point follows straightforwardly because $M$ and $\tilde{M}$ have the same partial solution functions if they are equivalent.

  The preservation of essential linearity follows because the coefficients of an affine mapping are identified if the mapping is known up to an $M$-null set.
\end{proof}

\begin{Eg}\label{eg:equivalent_scm}
  An iSCM with the single structural equation
  $$X = -X^3 + X + W$$
  where $X$ is endogenous, and $W$ is exogenous, is equivalent to the iSCM obtained by replacing the structural equation with
  $$X = \sqrt[3]{W}.$$
  Note that $X$ no longer appears on the r.h.s..
\end{Eg}

\begin{comment}
  This notion of equivalence is rather strong. 
\begin{Eg}\label{eg:self_cycle1}
  The iSCM with endogenous variable $X \in \RN$ and exogenous random variable $W \in \RN$ and structural equation
  $$X = W X + c$$
  is not equivalent to the iSCM obtained by replacing the structural equation with
  $$X = \begin{cases}
          \xi & W = 1 \\
          \frac{c}{1 - W} & W \ne 1
        \end{cases}$$
  no matter how we choose $\xi$, even when $P(W=1) = 0$. 
  If one does not model interventions on exogenous random variables, weaker notions of equivalence 
  can be used that replace the ``for all'' quantifier over values of exogenous random variables by the
  ``for almost all'' quantifier (see \cite{BFPM21}).
\end{Eg}
\end{comment}

We often make use of weaker notions of equivalence as well.
For simplicity of exposition, we provide the definitions only for simple iSCMs 
(the general definitions are provided in \cite{BFPM21} for SCMs).
%\footnote{\cite{BFPM21} also define a stronger notion of equivalent iSCMs which is useful when weakening the ``for all'' quantifier over values of exogenous random variables by the ``for almost all'' quantifier. Since here we will only work with the ``for all'' quantifier, there is no need to discuss that equivalence relation here.}
\begin{Def}\label{def:scm_obs_interv_equiv}
  Let $M = \lt \Sin, \Send, \Srnd, \CX, P, f \rt$ and
  $\tilde{M} = \ltuple \tilde{\Sin}, \tilde{\Send}, \tilde{\Srnd}, \tilde{\CX}, \tilde{P}, \tilde{f} \rtuple$ be two simple iSCMs
  and $O \subseteq \Send \cap \tilde{\Send}$ a subset.
  We say that:
  \begin{enumerate}
    \item $M$ and $\tilde{M}$ are \emph{observably equivalent w.r.t.\ $O$} if $\CX_O = \tilde{\CX}_O$, $\CX_{J \cap \tilde{J}} = \tilde{\CX}_{J \cap \tilde{J}}$ and their marginal Markov kernels coincide: 
      $$P_{M}(X_O \mid \doit(X_J)) = P_{\tilde{M}}(X_O \mid \doit(X_{\tilde J})).$$
    This has to be interpreted as both Markov kernels being a version of a Markov kernel $\CX_{J \cap \tilde{J}} \dshto \CX_O$,
    i.e., $P_{M}(X_O \mid \doit(X_J))$ must be constant in $x_{J \setminus \tilde{J}}$ and
    $P_{\tilde{M}}(X_O \mid \doit(X_{\tilde{J}}))$ must be constant in $x_{\tilde{J} \setminus J}$.\footnote{In other words,
      $$P_{M}(X_O \mid \doit(X_J)) = P_{M}(X_O \mid \doit(X_{J \cap \tilde{J}},\cancel{X_{J\sm \tilde{J}}})) =
      P_{\tilde{M}}(X_O \mid \doit(X_{J \cap \tilde{J}},\cancel{X_{\tilde{J}\sm J}})) = P_{\tilde{M}}(X_O \mid \doit(X_{\tilde{J}})).$$}
    \item $M$ and $\tilde{M}$ are \emph{interventionally equivalent w.r.t.\ $O$} if for every subset $T \subseteq O$ the intervened iSCMs $M_{\doit(T)}$ and $\tilde{M}_{\doit(T)}$ are observably equivalent w.r.t.\ $O \setminus T$.
  \end{enumerate}
  If $O = \Send \cap \tilde{\Send}$, we may omit the qualitifier ``w.r.t.\ $O$''.
  If $J = \tilde J =  \emptyset$, we also use the term \emph{observationally equivalent w.r.t.\ $O$} interchangebly with \emph{observably equivalent w.r.t.\ $O$}.
%  In case the variables in $O$ (and $J \cup \tilde{J}$) are considered observed, whereas the variables in 
%  $(V \cup W) \setminus O$ and $(\tilde{V} \cup \tilde{W}) \setminus O$ are considered latent, 
%  we sometimes omit the ``w.r.t.\ $O$'' in this terminology (i.e., we call $M$ and $\tilde{M}$ observably equivalent, interventionally equivalent, respectively).
\end{Def}
\begin{Rem}\label{rem:scm_interv_equiv}
  Note that $M$ and $\tilde{M}$ are interventionally equivalent w.r.t.\ $O \ins (V \cap \tilde{V})$ if and only if
  $\CX_O = \tilde{\CX}_O$, $\CX_{J \cap \tilde{J}} = \tilde{\CX}_{J \cap \tilde{J}}$
  and for all $T \ins O$:
    $$P_{M}(X_O \mid \doit(X_J,X_T)) = P_{\tilde{M}}(X_O \mid \doit(X_{\tilde J},X_T)).$$
\end{Rem}
More generally, one could define interventional equivalence not only with respect to an observed set of variables, but also with respect to a given set of interventions.

Equivalence is stronger than interventional equivalence, which in turn is stronger than observable equivalence.
\begin{Prp}\label{prp:pearls_ladder_levels_1_and_2}
  For simple iSCMs $M, \tilde{M}$ and a subset $O \subseteq \Send \cap \tilde{\Send}$:
  \begin{enumerate}
    \item If $M \equiv \tilde{M}$ then $M$ and $\tilde{M}$ are interventionally equivalent w.r.t.\ $O$.
    \item Interventional equivalence of $M$ and $\tilde{M}$ w.r.t.\ $O$ implies observable equivalence of $M$ and $\tilde{M}$ w.r.t.\ $O$.
  \end{enumerate}
\end{Prp}
\begin{proof}
  \begin{enumerate}
   \item
   Suppose $M \equiv \tilde{M}$. 
   Then for every $T \subseteq O$, $M_{\doit(T)} \equiv \tilde{M}_{\doit(T)}$.
   Since equivalent simple iSCMs have the same Markov kernels, $M_{\doit(T)}$ and $\tilde{M}_{\doit(T)}$ are observably equivalent w.r.t.\ $O \setminus T$ for every $T \subseteq O$.

  \item This is trivial (consider the intervention targeting $T = \emptyset$).
\end{enumerate}
\end{proof}
However, the reverse implication does not hold in general. This expresses that causal modeling is more refined than probabilistic modeling.
\begin{Eg}[Observable equivalence does not imply interventional equivalence]
Consider the SCM $M$ with
  \begin{align*}
    W &\sim \C{N}(\mu,\sigma^2), \\
    N &= W, \\
    C &= \alpha + \beta N,
  \end{align*}
and the SCM $\tilde{M}$ with
  \begin{align*}
    W &\sim \C{N}(\tilde{\mu},\tilde{\sigma}^2), \\
    C &= W, \\
    N &= \tilde{\alpha} + \tilde{\beta} C.
  \end{align*}
These SCMs are simple and their distributions are respectively 
  $$P_M(N,C) = \C{N}\left(\begin{pmatrix} \mu \\ \alpha + \beta \mu \end{pmatrix}, \begin{pmatrix} \sigma^2 & \beta\sigma^2 \\ \beta\sigma^2 & \beta^2\sigma^2\end{pmatrix}\right)$$
and
  $$P_{\tilde{M}}(N,C) = \C{N}\left(\begin{pmatrix} \tilde{\alpha} + \tilde{\beta} \tilde{\mu} \\ \tilde{\mu} \end{pmatrix}, \begin{pmatrix} \tilde{\beta}^2\tilde{\sigma}^2 & \tilde{\beta}\tilde{\sigma}^2 \\ \tilde{\beta}\tilde{\sigma}^2 & \tilde{\sigma}^2\end{pmatrix}\right)$$
  For certain parameter choices, they are observably equivalent (to be precise, they are observable equivalent iff 
  $\mu = \tilde{\alpha} + \tilde{\beta}\tilde{\mu}$ and $\tilde{\mu} = \alpha + \beta\mu$ and $\sigma^2 = \tilde{\beta}^2\tilde{\sigma}^2$ and $\beta \sigma^2 = \tilde{\beta}\tilde{\sigma}^2$ and $\beta^2 \sigma^2 = \tilde{\sigma}^2$).
  However, they are not interventionally equivalent except for very special parameter choices (to be precise, they are interventionally equivalent iff $\beta = \tilde{\beta} = 0$ and $\sigma^2=\tilde{\sigma}^2=0$ and $\mu = \tilde{\alpha}$ and $\tilde{\mu} = \alpha$).
\end{Eg}

We have seen that equivalence is preserved under various interventions. 
This is also the case for interventional equivalence (although one has to be careful with respect to which subset).
\begin{Prp}\label{prp:scm_interventions_preserve_interv_equiv}
  Assume that $M$ and $\tilde{M}$ are interventionally equivalent iSCMs w.r.t.\ $O \ins (V \cap \tilde{V})$. 
  Then, for $T \subseteq O \cup (J \cap \tilde{J})$:
    \begin{enumerate}
      \item $M_{\doit(T)}$ and $\tilde{M}_{\doit(T)}$ are interventionally equivalent w.r.t.\ $O \sm T$ (hard interventions with unspecified target value);
      \item $M_{\doit(X_T=\xi_T)}$ and $\tilde{M}_{\doit(X_T=\xi_T)}$ are interventionally equivalent w.r.t.\ $O \cup T$, for any $\xi_T \in \CX_T$ (hard intervention with specified target value);
      \item $M_{\doit(X_T\sim Q_T)}$ and $\tilde{M}_{\doit(X_T\sim Q_T)}$ are interventionally equivalent w.r.t.\ $O \cup T$, for any $Q_T \in \Pcal(\CX_T)$ (hard intervention with stochastic target value);
%      \item $M_{\doit(T\ot\hat{f}_T)} \equiv \tilde{M}_{\doit(T\ot\hat{f}_T)}$ for $T \subseteq V$ and $\hat f_T : \CX_J \times \CX_W \times \CX_V \to \CX_T$ a measurable function
%        (soft interventions);
      \item $M_{\doit(I_T)}$ and $\tilde{M}_{\doit(I_T)}$ are interventionally equivalent w.r.t.\ $O$ (adding intervention variables).
    \end{enumerate}
\end{Prp}
\begin{proof}
  \Joris{TODO: cleanup and final check}
  \begin{enumerate}
    \item This follows from the commutativity of hard interventions (Proposition~\ref{prp:interventions_commute}), plus some bookkeeping, as we show now in detail.
  Since $M$ and $\tilde{M}$ are interventionally equivalent w.r.t.\ $O$,
  for all subsets $S \ins O$, $M_{\doit(S)}$ and $\tilde{M}_{\doit(S)}$ are observably equivalent w.r.t.\ $O \sm S$.
  Let $T \subseteq O \cup (J \cap \tilde{J})$.
  $M_{\doit(T)}$ and $\tilde{M}_{\doit(T)}$ are interventionally equivalent w.r.t.\ $O\sm T$ if for all $\tilde{T} \ins O\sm T$,
  $(M_{\doit(T)})_{\doit(\tilde{T})}$ and $(\tilde{M}_{\doit(T)})_{\doit(\tilde{T})}$ are observably equivalent w.r.t.\ $(O\sm T) \sm \tilde{T}$.
  That is, if $M_{\doit(T \cup \tilde{T})}$ and $\tilde{M}_{\doit(T \cup \tilde{T})}$ are observably equivalent w.r.t.\ $O \sm (T \cup \tilde{T})$ for all $\tilde{T} \ins O\sm T$.
  That is, if $M_{\doit((T \cup \tilde{T}) \sm (J \cap \tilde{J}))}$ and $\tilde{M}_{\doit((T \cup \tilde{T}) \sm (J \cap \tilde{J}))}$ are observably equivalent w.r.t.\ $O \sm ((T \cup \tilde{T}) \sm (J \cap \tilde{J}))$ for all $\tilde{T} \ins O\sm T$.
  Note that for all $\tilde{T} \ins O\sm T$, $(T \cup \tilde{T}) \sm (J \cap \tilde{J}) \ins O$.
  \Claude{Does not fully close the final deduction.}
  \end{enumerate}

  For the next two cases, $\CX_O = \tilde{\CX}_O$ and $\CX_{J \cap \tilde{J}} = \tilde{\CX}_{J \cap \tilde{J}}$ holds since
  $M$ and $\tilde{M}$ are interventionally equivalent w.r.t.\ $O$. 
  This same property needs to hold for the intervened iSCMs $M_{\doit(T\dots)}$ and $\tilde{M}_{\doit(T\dots)}$.

\begin{enumerate}[resume*]
\item
  Let $T \subseteq O \cup (J \cap \tilde{J})$ and $\xi_T \in \CX_T$.
  To show: $M_{\doit(X_T=\xi_T)}$ and $\tilde{M}_{\doit(X_T=\xi_T)}$ are interventionally equivalent w.r.t.\ $O \cup T$, that is
  $\CX_O = \tilde{\CX}_O$, $\CX_{J \cap \tilde{J}} = \tilde{\CX}_{J \cap \tilde{J}}$, and for all $\tilde{T} \ins O \cup T$:
  $$P_{M_{\doit(X_T=\xi_T)}}(X_{O \cup T} \mid \doit(X_{J\sm T},X_{\tilde{T}})) = P_{\tilde{M}_{\doit(X_T=\xi_T)}}(X_{O\cup T} \mid \doit(X_{\tilde J \sm T},X_{\tilde{T}})).$$

  Note that $\doit(X_T=\xi_T)$ can be obtained by $\doit(T)$ followed by evaluating in $X_T=\xi_T$:
      $$P_{M_{\doit(X_T=\xi_T)}}(X_{O \cup T} \mid \doit(X_{J\sm T \cup \tilde{T}})) = P_M(X_{O \cup T} \mid \doit(X_{J \sm T \cup \tilde{T}}),\doit(X_T=\xi_T))$$
      $$P_{\tilde{M}_{\doit(X_T=\xi_T)}}(X_{O \cup T} \mid \doit(X_{\tilde{J}\sm T\cup \tilde{T}})) = P_{\tilde{M}}(X_{O \cup T} \mid \doit(X_{\tilde{J} \sm T \cup \tilde{T}}),\doit(X_T=\xi_T))$$
  So we want that
  $$P_M(X_{O} \mid \doit(X_{J \cup T \cup \tilde{T}})) = P_{\tilde{M}}(X_{O} \mid \doit(X_{\tilde{J} \cup T \cup \tilde{T}}))$$
  which we know holds for all $T \cup \tilde{T} \ins O \cup (J \cap \tilde{J})$ by the interventional equivalence of $M$ and $\tilde{M}$ w.r.t.\ $O$.

\item
  Let $T \subseteq O \cup (J \cap \tilde{J})$ and $Q_T \in \Pcal(\CX_T)$.
  We can use again that
  $$P_M(X_{O} \mid \doit(X_{J \cup T \cup \tilde{T}})) = P_{\tilde{M}}(X_{O} \mid \doit(X_{\tilde{J} \cup T \cup \tilde{T}}))$$
  for all $T \cup \tilde{T} \ins O \cup (J \cap \tilde{J})$ by the interventional equivalence of $M$ and $\tilde{M}$ w.r.t.\ $O$.
  Note that $\doit(X_T\sim Q_T)$ can be obtained by $\doit(T)$ followed by putting distribution $Q_T$ on $X_T$:
      $$P_{M_{\doit(X_T\sim Q_T)}}(X_{O \cup T} \mid \doit(X_{J\sm T \cup \tilde{T}})) = P_M(X_{O \cup T} \mid \doit(X_{J \sm T \cup \tilde{T}}),\doit(X_T)) \circ Q_T(X_T)$$
      $$P_{\tilde{M}_{\doit(X_T\sim Q_T)}}(X_{O \cup T} \mid \doit(X_{\tilde{J}\sm T\cup \tilde{T}})) = P_{\tilde{M}}(X_{O \cup T} \mid \doit(X_{\tilde{J} \sm T \cup \tilde{T}}),\doit(X_T)) \circ Q_T(X_T)$$
  Hence, for all $\tilde{T} \ins O \cup T$:
      $$P_{M_{\doit(X_T\sim Q_T)}}(X_{O \cup T} \mid \doit(X_{J\sm T}),\doit(X_{\tilde{T}})) = P_{\tilde{M}_{\doit(X_T\sim Q_T)}}(X_{O \cup T} \mid \doit(X_{\tilde{J}\sm T}),\doit(X_{\tilde{T}}))$$
  which implies that $M_{\doit(X_T\sim Q_T)}$ is interventionally equivalent to $\tilde{M}_{\doit(X_T\sim Q_T)}$ w.r.t.\ $O \cup T$.
  \item For $T \ins (J \cap \tilde{J})$, the operation of adding intervention variables $\doit(I_{T})$ only adds indices $I_j := j$ for $j \in T \cap J$.
    Therefore, we can consider $T \subseteq O$ without loss of generality, employing Proposition~\ref{prp:adding_intervention_variables_one_by_one}.

%  \begin{align*}
%    (M_{\doit(I_T)})_{\doit(\tilde{T})} 
%    &= (M_{\doit(I_T)})_{\doit((\tilde{T}\sm T) \cup (\tilde{T} \cap T))} \\
%    &= \big((M_{\doit(T)})_{\doit(\tilde{T}\sm T)}\big)_{\doit(\tilde{T} \cap T)} \\
%    &= \big((M_{\doit(\tilde{T}\sm T)})_{\doit(I_T)}\big)_{\doit(\tilde{T} \cap T)}
%  \end{align*}
%  and similarly for $\tilde{M}$.
  For $\tilde{T} \ins O$, we get:
  \begin{align*}
    & P_{(M_{\doit(I_T)})_{\doit(\tilde{T})}}(X_{V\sm\tilde{T}} \mid \doit(X_{J\cup\tilde{T}},X_{I_T}=x_{I_T})) \\
    {} = {} & P_{\big((M_{\doit(\tilde{T}\sm T)})_{\doit(I_T)}\big)_{\doit(\tilde{T} \cap T)}}(X_{V\sm\tilde{T}} \mid \doit(X_{J\cup\tilde{T}},X_{I_T}=x_{I_T})) \\
    {} = {} & P_{\big((M_{\doit(\tilde{T}\sm T)})_{\doit(I_{T\sm\tilde{T}})}\big)_{\doit(\tilde{T} \cap T)}}(X_{V\sm\tilde{T}} \mid \doit(X_{J\cup\tilde{T}},X_{I_{T\sm\tilde{T}}}=x_{I_{T\sm\tilde{T}}})) \\
    {} = {} & P_{(M_{\doit(\tilde{T})})_{\doit(I_{T\sm\tilde{T}})}}(X_{V\sm\tilde{T}} \mid \doit(X_{J\cup\tilde{T}},X_{I_{T\sm\tilde{T}}}=x_{I_{T\sm\tilde{T}}})) \\
%    {} = {} & P_{M_{\doit(\tilde{T})}}(X_{V\sm\tilde{T}} \mid \doit(X_{J\cup\tilde{T}},X_{T_2\sm\tilde{T}}=x_{I_{T_2\sm\tilde{T}}})) \\
    {} = {} & P_M(X_{V\sm\tilde{T}} \mid \doit(X_{J\cup\tilde{T}},X_{T_2\sm\tilde{T}}=x_{I_{T_2\sm\tilde{T}}})) 
  \end{align*}
  \Claude{$T_2$ undefined (occurs multiple times here), should be $T$.}
  Similar reasoning for $\tilde{M}$ gives:
  \begin{align*}
    & P_{(\tilde{M}_{\doit(I_T)})_{\doit(\tilde{T})}}(X_{\tilde{V}\sm\tilde{T}} \mid \doit(X_{\tilde{J}\cup\tilde{T}},X_{I_T}=x_{I_T})) \\
    {} = {} & P_{\tilde{M}}(X_{\tilde{V}\sm\tilde{T}} \mid \doit(X_{\tilde{J}\cup\tilde{T}},X_{T_2\sm\tilde{T}}=x_{I_{T_2\sm\tilde{T}}}))
  \end{align*}
  Since $M$ and $\tilde{M}$ are interventionally equivalent w.r.t.\ $O$, 
  \begin{align*}
    {}   {} & P_M(X_{O\sm\tilde{T}} \mid \doit(X_{J\cup\tilde{T}},X_{T_2\sm\tilde{T}}=x_{I_{T_2\sm\tilde{T}}})) \\
    {} = {} & P_{\tilde{M}}(X_{O\sm\tilde{T}} \mid \doit(X_{\tilde{J}\cup\tilde{T}},X_{T_2\sm\tilde{T}}=x_{I_{T_2\sm\tilde{T}}}))
  \end{align*}
  Combining this gives:
%  $$ P_{(M_{\doit(I_T)})_{\doit(\tilde{T})}}(X_{O\sm\tilde{T}} \mid \doit(X_{J\cup\tilde{T}},X_{I_T}=x_{I_T})) 
%     = P_{(\tilde{M}_{\doit(I_T)})_{\doit(\tilde{T})}}(X_{O\sm\tilde{T}} \mid \doit(X_{\tilde{J}\cup\tilde{T}},X_{I_T}=x_{I_T})) $$
%  To show:
  $$
   P_{(M_{\doit(I_T)})_{\doit(\tilde{T})}}(X_{O\sm\tilde{T}} \mid \doit(X_{J\cup\tilde{T}},X_{I_T})) 
   = 
   P_{(\tilde{M}_{\doit(I_T)})_{\doit(\tilde{T})}}(X_{O\sm\tilde{T}} \mid \doit(X_{\tilde{J}\cup\tilde{T}},X_{I_T})) $$
   Hence, $M_{\doit(I_T)}$ and $\tilde{M}_{\doit(I_T)}$ are interventionally equivalent w.r.t.\ $O$.
%  If $\tilde{T} \ins O \cup I_T$, note that
%  $$(M_{\doit(I_T)})_{\doit(\tilde{T})} = (M_{\doit(I_T)})_{\doit(\tilde{T}\sm I_T)}$$
%  and similarly for $\tilde{M}$. 
%  Since $(M_{\doit(I_T)})_{\doit(\tilde{T}\sm I_T)}$ is observationally equivalent w.r.t.\ $O \sm T$, we conclude that
%    $$P_{(M_{\doit(I_T)})_{\doit(\tilde{T})}}(X_{V\sm\tilde{T}} \mid \doit(X_{J\cup\tilde{T}},X_{I_T})) \\
%    = P_{(\tilde{M}_{\doit(I_T)})_{\doit(\tilde{T})}}(X_{\tilde{V}\sm\tilde{T}} \mid \doit(X_{\tilde{J}\cup\tilde{T}},X_{I_T}))$$
%  even for all $\tilde{T} \ins O \cup I_T$.
  \end{enumerate}
\end{proof}

\subsection{Marginalizations}

When modeling a system, we sometimes want to ``hide'' details of a subsystem. 
The following operation on iSCMs that we call ``marginalization'' is a causal analogue of the marginalization of probability distributions.
The computer program analogy of the marginalization operation is to hide details within a subroutine.
Intuitively, a marginalization of an iSCM over a subset of endogenous variables $L$ is obtained by first solving a subsystem (the structural equations corresponding to the endogenous variables in $L$) followed by substituting the solution function of the subsystem into the remaining structural equations (corresponding to the endogenous variables in $V\setminus L$).
\begin{Def}\label{def:scm_marginalization}
  Let $M = \lt \Sin, \Send, \Srnd, \CX, P, f \rt$ be an iSCM and $L \subseteq \Send$.
  Suppose that $M$ is essentially uniquely solvable w.r.t.\ $L$ and let $g^{[L]} : \CX_J \times \CX_{(V\sm L)} \times \CX_W \to \CX_L$ be a corresponding partial solution function of $M$ w.r.t.\ $L$.
%  that satisfies
%    $$x_L = f_L(x) \iff x_L = g^{[L]}(x_J,x_{(V\sm L)},x_W)\quad M\as.$$
  Then we call $M_{\setminus L} = \ltuple \Sin, V \sm L, \Srnd, \CX_{\Sin} \times \CX_{V \sm L} \times \CX_{\Srnd}, P, f^{\sm L} \rtuple$ with
  $$f^{\sm L}(\xJ,x_{V \sm L},\xW) := f_{V \sm L}(\xJ,x_{V \sm L},g^{[L]}(\xJ,x_{V \sm L},\xW),\xW)$$
  a \emph{marginalization of $M$ over $L$}.
\end{Def}
% \begin{Def}\label{def:scm_marginalization}
%  Let $M = \lt \Sin, \Send, \Srnd, \CX, P, f \rt$ be an iSCM and $L \subseteq \Send$ such that the submodel $M^{[L]}$ is essentially uniquely solvable.
%  Let $g^{[L]} : \CX_{\Sin} \times \CX_{V \sm L} \times \CXW \to \CX_L$ be an (essentially unique) solution function for $M^{[L]}$.
%  Then we call $M_{\setminus L} = \ltuple \Sin, V \sm L, \Srnd, \CX_{\Sin} \times \CX_{V \sm L} \times \CX_{\Srnd}, P, \tilde{f} \rtuple$ with
%  $$\tilde{f}(\xJ,x_{V \sm L},\xW) = f_{V \sm L}(\xJ,x_{V \sm L},g^{[L]}(\xJ,x_{V \sm L},\xW),\xW)$$
%  a \emph{marginalization of $M$ over $L$}.
%\end{Def}
If $g^{[L]}$ and $\tilde g^{[L]}$ are two partial solution functions of $M$ w.r.t.\ $L$, then their corresponding marginalizations of $M$ may not be identical, but they are equivalent. Even more:
\begin{Prp}\label{prp:scm_marginalization_equivalence}
  If $M \equiv \tilde{M}$ are equivalent iSCMs, $g^{[L]}$ is a partial solution function of $M$ w.r.t.\ $L$, and $\tilde g^{[L]}$ is a partial solution function of $\tilde{M}$ w.r.t.\ $L$, then the marginalizations $M_{\sm L}$ (corresponding to $g^{[L}]$) and $\tilde{M}_{\sm L}$ (corresponding to $\tilde g^{[L]}$) are equivalent.
\end{Prp}
\begin{proof}
  Since $\tilde{g}^{[L]}$ is a partial solution function of $\tilde{M}$ w.r.t.\ $L$, and $M \equiv \tilde{M}$, $\tilde{g}^{[L]}$ is also a partial solution function of $M$ w.r.t.\ $L$.
  Since $M$ is essentially uniquely solvable w.r.t.\ $L$, $g^{[L]} = \tilde{g}^{[L]}\ M\as$.
  Therefore, for $v \in V\sm L$:
  $$[x_v = f_v(x_J,x_{V\sm L},g^{[L]}(x),x_W) \iff x_v = \tilde{f}_v(x_J,x_{V\sm L},\tilde{g}^{[L]}(x),x_W)] \qquad M\as$$
  This implies, for $v \in V\sm L$:
  $$[x_v = f^{\sm L}_v(x) \iff x_v = \tilde{f}^{\sm L}_v(x)] \qquad M_{\sm L}\as.$$
  Indeed, if a property $\pi(x_J,x_W,x_{V\sm L})$ holds $M\as$, then it holds $M_{\sm L}\as$.
  %Indeed, if something holds `for all $x_J$, for $P(X_W)$-almost all $x_W$, for all $x_V$', then it holds `for all $x_J$, for $P(X_W)$-almost all $x_W$, for all $x_{V\sm L}$'.
\end{proof}
This shows that even though marginalizations are not unique at the iSCM level, they are unique at the level of equivalence classes of iSCMs.

For simple iSCMs, marginalizations are obviously defined over all subsets $L \subseteq \Send$.
\begin{Not}\label{not:scm_marginalization}
  We will also use the notation $M_O := M_{\sm (V \sm O)}$ for a subset $O \subseteq V$ of endogenous variables
  when we want to emphasize which endogenous variables remain after marginalization.
\end{Not}
%Note that all marginalizations of $M$ are equivalent, i.e., the marginalized causal mechanism is essentially unique.

Marginalization preserves (partial) essentially unique solvability.
\begin{Lem}\label{lem:marginalization_preserves_unique_solvability}
  Let $M = \lt \Sin, \Send, \Srnd, \CX, P, f \rt$ be an iSCM and $L \subseteq \Send$ such that $M$ is essentially uniquely solvable w.r.t.\ $L$.
  Suppose $M$ is also essentially uniquely solvable w.r.t.\ $O \ins V$ with $L \ins O$, and partial solution function $g^{[O]}$. 
  Then its marginalization $\tilde M := M_{\setminus L}$ is essentially uniquely solvable w.r.t.\ $O \sm L$, with partial solution function $\tilde{g}^{[O\sm L]} = (g^{[O]})_{O\sm L}$.
  % that is, the canonical projection onto $O \sm L$ of $g^{[O]}$.
\end{Lem}
\begin{proof}
  Denote by $f^{\sm L}$ the causal mechanism of the marginalization $M_{\sm L}$ constructed from the partial solution function $g^{[L]} : \CX_{\Sin} \times \CX_{(V \sm L)} \times \CXW \to \CX_L$.
  From the essentially unique solvability of $M$ w.r.t.\ $O$, and the essentially unique solvability of $M$ w.r.t.\ $L$, we derive that
%  $M$ is uniquely solvable means there exists a measurable function $g : \CX_{J \cup W} \to \CX_V$ such  that
%  for all $x \in \CX$,
  $M\as$:
  \begin{equation*}\begin{split}
  & \begin{cases}
    x_L = g^{[O]}_L(x_J,x_{V\sm O},x_W) \\
    x_{O\sm L} = g^{[O]}_{O\sm L}(x_J,x_{V\sm O},x_W) 
  \end{cases}
  \iff x_O = g^{[O]}(x_J,x_{V\sm O},x_W) 
  \iff x_O = f_O(x) \\
  & \iff \begin{cases}
    x_L = f_L(x) \\
    x_{O \sm L} = f_{O \sm L}(x) 
  \end{cases} 
  \iff \begin{cases}
    x_L = g^{[L]}(x_J,x_{V\sm L},x_W) \\
    x_{O \sm L} = f_{O \sm L}(x_J,x_{V\sm L},x_L,x_W) 
  \end{cases} \\
  & \iff \begin{cases}
    x_L = g^{[L]}(x_J,x_{V\sm L},x_W) \\
    x_{O\sm L} = f_{O \sm L}(x_J,x_{V\sm L},g^{[L]}(x_J,x_{V\sm L},x_W),x_W) 
  \end{cases} \\
  & \iff \begin{cases}
    x_L = g^{[L]}(x_J,x_{V\sm L},x_W) \\
    x_{O \sm L} = f_{O \sm L}^{\sm L}(x_J,x_{V\sm L},x_W).
  \end{cases}
  \end{split}\end{equation*}
Hence:
  $$x_{O \sm L} = g^{[O]}_{O \sm L}(x_J,x_{V\sm O},x_W) \iff x_{O \sm L} = f_{O \sm L}^{\sm L}(x_J,x_{V\sm L},x_W)\quad M\as.$$
  A property $\pi(x_J,x_W,x_{V\sm L})$ that holds $M\as$, also holds $M_{\sm L}\as$, hence:
  $$x_{O \sm L} = g^{[O]}_{O \sm L}(x_J,x_{V\sm O},x_W) \iff x_{O \sm L} = f_{O \sm L}^{\sm L}(x_J,x_{V\sm L},x_W)\quad M^{\sm L}\as.$$
  Therefore, the marginalized iSCM $M_{\setminus L}$ is essentially uniquely solvable w.r.t.\ $O\sm L$, with essentially unique partial solution function $g^{[O]}_{O \sm L} = \pr_{O\sm L} \circ g^{[O]}$,
  where $\pr_{O\sm L} : \CXV \to \CX_{O\sm L} : x \mapsto x_{O\sm L}$ is the canonical projection on $O \sm L$.\footnote{The converse does not necessarily hold: even if the marginalized iSCM $M_{\sm L}$ is essentially uniquely solvable w.r.t.\ $O \sm L$, the iSCM $M$ may not be essentially uniquely solvable w.r.t.\ $O$: $M^{\sm L}\as$ does not imply $M\as$. \Joris{check!}}
\end{proof}
\begin{Rem}
This also directly implies that if $M$ is essentially uniquely solvable and its marginalization $M_{\setminus L}$ over $L \subseteq V$ is defined, the Markov kernel of the marginalization is obtained by marginalizing the original Markov kernel:
$$P_{M_{\setminus L}}(X_{V\setminus L} \mid \doit(X_J)) = P_{M}(X_{V\setminus L} \mid \doit(X_J)).$$
This explains the name `marginalization'.
\end{Rem}
It is now easy to see that simplicity is preserved under marginalization.
\begin{Prp}\label{prp:marginalization_preserves_simplicity}
  Let $M = \lt \Sin, \Send, \Srnd, \CX, P, f \rt$ be a simple iSCM. 
  For any $L \subseteq \Send$, its marginalization $M_{\setminus L}$ is also simple.
\end{Prp}
\begin{proof}
  Let $\tilde T \ins V \sm L$.
  From Lemma~\ref{lem:marginalization_preserves_unique_solvability} it follows that since $M$ is essentially uniquely solvable w.r.t.\ $\tilde{T} \cup L$, $M_{\sm L}$ is essentially uniquely solvable w.r.t.\ $\tilde{T}$.
  Hence $M_{\sm L}$ is simple.
\end{proof}

Under certain conditions, interventions and marginalization commute (i.e., it does not matter in which order we apply them).
\begin{Prp}\label{prp:interventions_marginalizations_commute}
Let $M = \lt \Sin, \Send, \Srnd, \CX, P, f \rt$ be an iSCM.
    For $L \subseteq V$ such that $M$ is essentially uniquely solvable w.r.t.\ $L$, and an intervention target $T \ins V \cup J$ such that $L \cap T = \emptyset$, an intervention targeting $T$ (hard, soft, or adding intervention variables) commutes with marginalizing out $L$.
%    $$(M_{\doit(T\dots)})_{\setminus L} \equiv (M_{\setminus L})_{\doit(T\dots)}.$$
%    Explicitly:
    \begin{enumerate}
      \item $(M_{\doit(T\dots)})_{\sm L} \equiv (M_{\sm L})_{\doit(T\dots)}$ for $T \subseteq (V \sm L) \cup J$ (hard interventions of all three types);
      \item $(M_{\doit(T\ot\hat{f}_T)})_{\sm L} \equiv (M_{\sm L})_{\doit(T\ot\hat{f}_T)}$ for $T \subseteq V \sm L$ and $\hat f_T : \CX_J \times \CX_W \times \CX_V \to \CX_T$ a measurable function
        (soft interventions); \Joris{We could more generally also define soft interventions on input nodes!}
      \item $(M_{\doit(I_T)})_{\sm L} \equiv (M_{\sm L})_{\doit(I_T)}$ for $T \subseteq (V \sm L) \cup J$ (adding intervention variables);
    \end{enumerate}
\end{Prp}
\begin{proof}
  Note that each intervention preserves the causal mechanism components $f_{V\sm L}$ of the variables not targeted by the intervention. 
  Hence, we may use the same partial solution function w.r.t.\ $L$ for both $M$ and $M_{\doit(\dots)}$ to construct their respective marginalizations.
  This even gives identities.
  If different partial solution functions are used in the construction of the marginalizations, then we still obtain equivalence.
\end{proof}

We will show that for simple iSCMs, the marginalization operation preserves the causal semantics on the remaining variables. 
A key step is the following proposition, which gives conditions under which it does not matter whether we marginalize at once or in steps.
\begin{Prp}\label{prp:marginalization_commutes}
  Let $M = \lt \Sin, \Send, \Srnd, \CX, P, f \rt$ be an iSCM and $L_1, L_2 \subseteq \Send$ such that $L_1 \cap L_2 = \emptyset$.
  If $M$ is essentially uniquely solvable w.r.t.\ $L_1$, and $(M_{\setminus L_1})$ is essentially uniquely solvable w.r.t.\ $L_2$, then $M$ is essentially uniquely solvable w.r.t.\ $L_1 \cup L_2$, and in that case it does not matter if we first marginalize over $L_1$ and then $L_2$, or both at once, i.e:
  $$(M_{\setminus L_1})_{\setminus L_2} \equiv M_{\setminus (L_1\cup L_2)}.$$
\end{Prp}
\begin{proof}
  Write $K_1 = V \setminus L_1$.
  Let $g^{[L_1]} : \CX_{J\cup K_1} \times \CXW \to \CX_{L_1}$ be the essentially unique solution function of $M$ w.r.t.\ $L_1$:
  $$x_{L_1} = g^{[L_1]}(x_J,x_{K_1},x_W) \iff x_{L_1} = f_{L_1}(x) \quad M\as$$
  Let $\tilde{f} : \CX_J \times \CX_{K_1} \times \CXW \to \CX_{K_1}$ with
  $$\tilde{f}(x_J,x_{K_1},x_W) = f_{K_1}(x_J,x_{K_1},g^{[L_1]}(x_J,x_{K_1},x_W),x_W)$$
  be the causal mechanism of the marginal iSCM $M_{\setminus L_1}$.

  If $M_{\setminus L_1}$ is uniquely solvable w.r.t.\ $L_2$, it has an essentially unique solution function
  $\tilde{g}^{[L_2]} : \CXJ \times \CX_{V\setminus (L_1\cup L_2)} \times \CXW \to \CX_{L_2}$:
  $$x_{L_2} = \tilde{g}^{[L_2]}(x_J,x_{V\setminus(L_1\cup L_2)},x_W) \iff x_{L_2} = \tilde{f}_{L_2}(x_J,x_{L_2},x_{K_1\setminus L_2},x_W) \quad M_{\sm L_1}\as$$
  For a property $\pi(x_J,x_{V\sm L_1},x_W)$: $\pi$ holds $M\as$ if and only if $\pi$ holds $M_{\sm L_1}\as$, hence
  $$x_{L_2} = \tilde{g}^{[L_2]}(x_J,x_{V\setminus(L_1\cup L_2)},x_W) \iff x_{L_2} = f_{L_2}(x_J,x_{K_1},g^{[L_1]}(x_J,x_{K_1},x_W),x_W) \quad M\as$$
  Define the function $h : \CXJ \times \CX_{V\setminus(L_1\cup L_2)} \times \CXW \to \CX_{L_1 \cup L_2}$ by
  \begin{align*}
    h_{L_1}(x_J,x_{V\setminus(L_1\cup L_2)},x_W) &= g^{[L_1]}(x_J,x_{K_1\setminus L_2},\tilde{g}^{[L_2]}(x_J,x_{V\setminus(L_1 \cup L_2)},x_W),x_W) \\
    h_{L_2}(x_J,x_{V\setminus(L_1\cup L_2)},x_W) &= \tilde{g}^{[L_2]}(x_J,x_{V\setminus(L_1 \cup L_2)},x_W)
  \end{align*}
  Then, $M\as$:
  \begin{equation*}\begin{split}
%    & x_{L_1\cup L_2} = f_{L_1 \cup L_2}(x) \\
%    \iff &
    & \begin{cases}
      x_{L_1} &= f_{L_1}(x) \\
      x_{L_2} &= f_{L_2}(x) 
    \end{cases} \\
    \iff &
    \begin{cases}
      x_{L_1} &= g^{[L_1]}(x_J,x_{K_1},x_W) \\
      x_{L_2} &= f_{L_2}(x_J,x_{K_1},x_{L_1},x_W) 
    \end{cases} \\
    \iff &
    \begin{cases}
      x_{L_1} &= g^{[L_1]}(x_J,x_{K_1},x_W) \\
      x_{L_2} &= f_{L_2}(x_J,x_{K_1},g^{[L_1]}(x_J,x_{K_1},x_W),x_W) 
    \end{cases} \\
    \iff &
    \begin{cases}
      x_{L_1} &= g^{[L_1]}(x_J,x_{K_1\setminus L_2},x_{L_2},x_W) \\
      x_{L_2} &= \tilde{g}^{[L_2]}(x_J,x_{V\setminus(L_1 \cup L_2)},x_W)
    \end{cases} \\
    \iff &
    \begin{cases}
      x_{L_1} &= g^{[L_1]}(x_J,x_{K_1\setminus L_2},\tilde{g}^{[L_2]}(x_J,x_{V\setminus(L_1 \cup L_2)},x_W),x_W) \\
      x_{L_2} &= \tilde{g}^{[L_2]}(x_J,x_{V\setminus(L_1 \cup L_2)},x_W)
    \end{cases} \\
    \iff &
      x_{L_1\cup L_2} = h(x_J,x_{V\setminus(L_1\cup L_2)},x_W) 
  \end{split}\end{equation*}
  where in the first equivalence we used the essentially unique solvability of $M$ w.r.t.\ $L_1$, 
  in the second equivalence we used substitution, 
  in the third equivalence we used the essentially unique solvability of $M_{\setminus L_1}$ w.r.t.\ $[L_2]$, 
  in the fourth equivalence we used substitution again, 
  and in the fifth equivalence we used the definition of $h$.
  Hence,
    $$x_{L_1 \cup L_2} = f_{L_1 \cup L_2}(x) \iff x_{L_1\cup L_2} = h(x_J,x_{V\setminus(L_1\cup L_2)},x_W) \qquad M\as$$
  Therefore, $h$ is the essentially unique solution function for $M$ w.r.t.\ $L_1\cup L_2$, which must therefore be essentially uniquely solvable w.r.t.\ $L_1 \cup L_2$. 
  By checking the definition, one sees that $(M_{\setminus L_1})_{\setminus L_2} = M_{\setminus(L_1 \cup L_2)}$.
  Because of the ambiguity in the choice of the partial solution functions, we conclude $(M_{\setminus L_1})_{\setminus L_2} \equiv M_{\setminus(L_1 \cup L_2)}$.
\end{proof}
\begin{Cor}\label{cor:marginalization_commutes_simple_scm}
  Let $M = \lt \Sin, \Send, \Srnd, \CX, P, f \rt$ be a simple iSCM.
  For $L_1, L_2 \subseteq \Send$ with $L_1 \cap L_2 = \emptyset$,
    $$(M_{\setminus L_1})_{\setminus L_2} \equiv (M_{\setminus L_2})_{\setminus L_1} \equiv M_{\setminus (L_1 \cup L_2)}.$$
\end{Cor}

These commutation relations and compatibilities now allow us to give a straightforward proof that
the causal semantics are preserved under marginalization.
While this holds generally for SCMs \cite{BFPM21}, we will here prove this for simple iSCMs.
\begin{Thm}\label{thm:various_equiv_marginalization}
  Let $M = \lt \Sin, \Send, \Srnd, \CX, P, f \rt$ be a simple iSCM, $L \subseteq \Send$, and
  $M_{\setminus L}$ its marginalization over $L$. Then $M$ and $M_{\setminus L}$ are observably
  and interventionally equivalent w.r.t.\ $V \setminus L$.
  % counterfactual equivalence moved to a corollary later on
\end{Thm}
\begin{proof}
  %The statement that $M$ and $M_O$ are observably equivalent w.r.t.\ $O$ is Lemma~\ref{lem:obs_equiv_marginalization}.
  Write $O = V \setminus L$.
  We first show that the marginal Markov kernels $P_{M}(X_O \mid \doit(X_J))$ and $P_{M_{\setminus L}}(X_O \mid \doit(X_J))$ are the same.
  The former is obtained as:
  $$P_{M}(X_O \mid \doit(X_J)) = (\pr_{O} \circ g)_*(P) = (g_{O})_*(P),$$
  where $g : \CXJ \times \CXW \to \CXV$ is the essentially unique solution function of $M$ and
  $\pr_{O} : \CX_{V \cup W} \to \CX_{O}$ is the canonical projection on the $O$ components. 
  The latter is obtained as:
  $$P_{M_{\setminus L}}(X_O \mid \doit(X_J)) = (g_O)_*(P)$$
  since by Lemma~\ref{lem:marginalization_preserves_unique_solvability}, $g_O$ is the (essentially unique) solution function of $M_{\setminus L}$.
  This means that both push-forwards are identical.

  Let $T \subseteq O$. Then $(M_{\setminus L})_{\doit(T)} \equiv (M_{\doit(T)})_{\setminus L}$ by
  Proposition~\ref{prp:interventions_marginalizations_commute}. 
  The observable equivalence of $M_{\doit(T)}$ and $(M_{\doit(T)})_{\setminus L}$ w.r.t.\ $O \setminus T$ hence implies the observable equivalence of $M_{\doit(T)}$ and $(M_{\setminus L})_{\doit(T)}$ w.r.t.\ $O \setminus T$ (using Proposition~\ref{prp:pearls_ladder_levels_1_and_2}).
  Since this holds for all $T \subseteq O$, $M$ and $M_{\setminus L}$ are interventionally equivalent w.r.t.\ $O$.
%  Finally, $(M_{\setminus L})^\twin = (M^\twin)_{\setminus (L \cup L')}$ by Proposition~\ref{prp:twinning_marginalization_commute}. Since $M^\twin$ and its marginalization
%  $(M^\twin)_{\setminus (L \cup L')}$ are interventionally equivalent w.r.t.\ $M \cup M' \cup W$, 
%  $M$ and $M_{\setminus L}$ are counterfactually equivalent w.r.t.\ $M \cup W$.
\end{proof}
So the marginalization operation indeed effectively hides the details of a subsystem, while preserving the causal semantics on the remaining part.
